\section{Introduction}
\label{sec:introduction}

% 中文: LLM 处理代码时表现出一种令人困惑的异质性：同一个模型可以轻松追踪一条变量赋值链，
% 却在加入算术后变得脆弱；能处理展开的直线计算，却在等价的循环语义上挣扎；对嵌套条件分支
% 自信地给出错误答案，却在简单 boolean short-circuit 上迟迟无法收敛。这种异质性不仅出现在
% code-specialized 模型中，也出现在 general-purpose LLM 处理代码时。单看任务准确率无法
% 解释这些差异的来源——两个 accuracy 相近的模型可能在内部以完全不同的方式失败。要理解
% "代码理解"究竟是一种能力还是多种，我们需要打开模型内部，观察代码推理在逐层计算中如何展开。
LLMs exhibit a puzzling heterogeneity when processing code: a single model can effortlessly trace a chain of variable assignments yet become brittle once arithmetic is introduced; it can handle unrolled straight-line computations yet struggle with semantically equivalent loops; it confidently produces wrong answers on nested conditionals yet fails to converge on simple boolean short-circuit evaluation. This heterogeneity appears not only in code-specialized models but also when general-purpose LLMs process code \citep{liu2024codemind, gu2024cruxeval, abdollahi2025demystifying, chen2025reval}. Task-level accuracy alone cannot explain the source of these differences---two models with similar accuracy may fail internally in entirely different ways. To understand whether ``code understanding'' is a single capability or many, we need to look inside the model and observe how code reasoning unfolds across layers of computation.

% 中文: 现有可解释性工作大多回答的是一个表征性问题：模型在某一层编码了什么信息。但对代码推理
% 而言，更关键的问题是：在每一层，模型是否已经解出了这个问题？而且"解出"本身还有程度之分——
% 答案信息可能已经存在于 hidden state 中，但尚未被组织成模型自己能够稳定利用的形式。区分
% "信息存在"与"信息可用"，正是理解不同 code primitive 为何引发不同内部失败模式的入口。
Most existing interpretability work addresses a representational question: what information is \textit{encoded} at a given layer \citep{belinkov2017neural, cunningham2023sparse, nostalgebraist2020, belrose2023eliciting}. For code reasoning, however, the more critical question is: \textbf{at each layer, has the model already \textit{solved} the problem?} Moreover, ``solved'' itself admits degrees---answer information may already be present in the hidden state yet not organized into a form that the model can stably utilize on its own. Distinguishing \textit{information availability} from \textit{information readiness} is precisely the entry point for understanding why different code primitives trigger different internal failure modes.

% 中文: 为研究这一问题，我们构建了一个面向 interpretability 的 code reasoning benchmark，
% 包含六类 synthetic 任务，沿 data flow 到 control flow 的谱系设计。每个实例都被约束为
% 单 digit 输出，以避免 multi-token decoding 干扰逐层分析。在此基础上，我们引入一个
% dual diagnostic framework。逐层 linear probe 判断正确答案是否已在 hidden state 中
% 线性可读（information availability）；基于 Patchscopes 的 Context-Stripped Decoding
% (CSD) 则测试模型能否在移除原始代码上下文后，仅凭同一个 hidden state 自行解码出答案
% （information readiness）。
To investigate this question, we construct an interpretability-oriented code reasoning benchmark comprising six families of synthetic tasks, designed along a spectrum from data flow to control flow. Every instance is constrained to a single-digit output, preventing multi-token decoding from confounding layer-wise analysis. On top of this benchmark, we introduce a dual diagnostic framework. Layer-wise linear probing determines whether the correct answer is already linearly readable in the hidden state (\textit{information availability}); Context-Stripped Decoding (CSD), built upon Patchscopes \citep{ghandeharioun2024patchscopes}, tests whether the model can decode the answer from the same hidden state on its own once the original code context is removed (\textit{information readiness}).

% 中文: 这一视角揭示了代码推理的一个结构化内部过程。跨任务地看，模型都会先经历一个 brewing
% 过程：答案信息先对外部 readout 可见，但还没有变成模型自己可稳定解码的形式。随后轨迹分叉
% 为四种 resolution outcome——Resolved、Overprocessed、Misresolved 和 Unresolved——
% 共覆盖 [XX]% 的样本。对 code understanding 最关键的发现是：不同代码原语会触发显著不同的
% outcome fingerprint 和 brewing dynamics。例如，计算上完全等价的 Loop 与 Loop-unrolled
% 在 outcome fingerprint 上显著分离，直接表明模型敏感的不只是底层运算，还包括 code primitive
% 本身；而跨规模比较则揭示出 brewing scaffold 在不同模型间保持稳定，变化的只是 resolution
% 的成功率。代码推理并不是单一能力，而是一组共享早期结构、但在 resolution 路径上显著分化的
% 内部子过程。
This perspective reveals a structured internal process underlying code reasoning. Across tasks, models first undergo a \textbf{brewing} process in which the answer becomes visible to an external readout before it is organized into a form the model itself can stably decode. Trajectories then diverge into four \textbf{resolution outcomes}---\textbf{Resolved}, \textbf{Overprocessed}, \textbf{Misresolved}, and \textbf{Unresolved}---which together cover [XX]\% of observed samples. The most consequential finding for code understanding is that different code primitives trigger markedly different outcome fingerprints and brewing dynamics. For instance, Loop and its computationally equivalent Loop-unrolled variant exhibit significantly separated outcome fingerprints, directly demonstrating that the model is sensitive not merely to the underlying computation but to the code primitive itself. Cross-scale comparisons further reveal that the brewing scaffold remains stable across models and scales; what changes is only the success rate of resolution. Code reasoning is not a monolithic capability but a collection of internal sub-processes that share an early structure yet diverge significantly along their resolution paths.

% 中文: 我们的贡献有三点：
% - 我们构建了一个用于逐层研究代码推理的 benchmark + dual diagnostic framework，将
%   probing 与 Context-Stripped Decoding 作为 information availability 与 information
%   readiness 的互补视角。
% - 我们刻画了一个 brewing-to-resolution structure：答案往往先变得外部可读，再变得对模型
%   自身可解码；其后轨迹分化为一个经因果支持的四路 resolution taxonomy。
% - 我们展示了 code primitive 层面的 mechanistic differences：不同代码原语诱导出不同的
%   outcome fingerprint 和 brewing dynamics，暴露出不同的 resolution bottleneck；同时
%   brewing scaffold 跨模型与跨规模保持稳定，而 outcome 分布随能力变化，表明 mechanism
%   与 capability 可分离。
Our contributions are threefold:
\begin{itemize}
    \item We construct a \textbf{benchmark paired with a dual diagnostic framework} for layer-wise study of code reasoning, combining probing and Context-Stripped Decoding as complementary lenses on information availability and information readiness.
    \item We characterize a \textbf{brewing-to-resolution structure}: the answer typically becomes externally readable before it becomes self-decodable by the model; subsequent trajectories diverge into a causally supported four-way resolution taxonomy.
    \item We demonstrate \textbf{code-primitive-level mechanistic differences}: different code primitives induce distinct outcome fingerprints and brewing dynamics, exposing different resolution bottlenecks; meanwhile, the brewing scaffold remains stable across models and scales while outcome distributions shift with capability, indicating that mechanism and capability are separable.
\end{itemize}
